% 本文件由 LaTeX 排版 Agent 根据有效 translation.json 自动生成。
% author=Councils; work_id=3806; title=劳狄刻雅会议（第四世纪）

\chapter{劳狄刻雅会议（第四世纪）}

% structure: p0001 source=h1 target=omitted reason=page-title-already-rendered-by-parent

% structure: p0002 source=h2 target=section reason=relative-heading-rank; base=h2

\section{历史引言}

会议所在的劳狄刻雅是弗里吉亚·帕卡提亚纳的劳狄刻雅，又称利库斯河畔劳狄刻雅，需谨慎区分于叙利亚的劳狄刻雅。这一点是确定的，但关于会议的确切日期则多有讨论。\Person{彼得·德·马尔卡}将其定于365年，但\Person{帕吉}在其对\Person{巴罗尼乌斯}《编年史》的《批评》中似乎推翻了德·马尔卡所依据的论点，并与\Person{戈托弗雷德}一致将其定于约363年。初看之下，似乎第七条诫命给出了线索可确定日期，因为其中提到了福提努斯派，而福提努斯主教在第四世纪中期开始显赫，并于344年在安提约基雅被优西比乌派诅咒，于345年在米兰被正统派诅咒；最后，在数次其他谴责后，他于366年死于流放。但\Quote{福提努斯派}一词是否为插入语尚不完全确定。关于日期，也许可以从描述弗里吉亚的词语\OriginalTerm{帕卡提亚纳的}{\GreekText{Πακατιανῆς}}中得出一些信息，因为这一划分很可能在343年萨尔迪卡会议时尚未确立。\Person{赫费勒}得出结论：\Quote{在这种情况下，最好与\Person{雷米·塞利耶}、\Person{蒂勒蒙}等人一同，将劳狄刻雅会议的召开时间大致定在343年至381年之间，即萨尔迪卡会议与第二次大公会议之间，并放弃寻找更精确日期的尝试。}\par

但由于这次会议的教规传统上排在安提约基雅会议之后、第一次君士坦丁堡会议之前，我遵循了这一顺序。\Person{赫费勒}说，\Quote{许多始于第六世纪乃至第五世纪的古老教规汇编}正是如此排列。诚然，\Person{玛窦·布拉斯塔雷斯}将这些教规置于萨尔迪卡会议之后，但根据《教规大全》，昆尼塞克特会议在其第二条教规中，以及\Person{教宗良四世}，都赋予了它们在本卷中的位置。\par

% structure: p0005 source=h2 target=section reason=relative-heading-rank; base=h2

\section{教规}

\Emph{关于在弗里吉亚·帕卡提亚纳的劳狄刻雅城举行的会议教规，该会议由来自亚洲各省的许多可敬教父共同召开。}\par

在弗里吉亚·帕卡提亚纳的劳狄刻雅聚集的圣会议，来自亚洲各地，制定了以下教会法规。\par

% structure: p0008 source=h3 target=subsection reason=relative-heading-rank; base=h2

\subsection{第一条教规}

依照教会教规，对于自由合法地进入第二次婚姻、且之前未有秘密婚姻者，应在他们经过一段时间的祈祷和守斋后，恩准其领受共融。\par

% structure: p0010 source=h3 target=subsection reason=relative-heading-rank; base=h2

\subsection{第二条教规}

凡在各种事上犯罪者，若坚持告解和补赎的祈祷，并完全脱离其过犯，将在根据其罪行性质所指定的补赎时期后，因天主的慈悲和良善，重新获准领受共融。\par

% structure: p0012 source=h3 target=subsection reason=relative-heading-rank; base=h2

\subsection{第三条教规}

新受洗者不应被擢升至司铎职。\par

% structure: p0014 source=h3 target=subsection reason=relative-heading-rank; base=h2

\subsection{第四条教规}

司铎职之人不应放贷收取利息，也不应收取所谓的\Quote{半倍利}（\OriginalTerm{半倍利}{hemioliæ}）。\par

\EditorialNote{狄奥尼修斯·埃克西古斯和伊西多尔将此条列为第五条，而我们将第五条列为第四条。}\par

% structure: p0017 source=h3 target=subsection reason=relative-heading-rank; base=h2

\subsection{第五条教规}

\Emph{祝圣}不应在有听众在场时举行。\par

\Person{巴尔萨蒙}：\Quote{此教规将选举称为\Quote{按手}，并说，由于在选举中常对被选举者说不恰当的话，因此不应在任何可能来旁听的人面前举行。}\par

\Person{佐纳拉斯}也同意此处指选举，但\Person{阿里斯特努斯}不同意，并将之理解为严格意义上的祝圣，如下：祝圣祷文不应大声诵读以免被民众听到。\par

% structure: p0021 source=h3 target=subsection reason=relative-heading-rank; base=h2

\subsection{第六条教规}

异端者在其仍坚持异端期间，不得进入天主的殿宇。\par

% structure: p0023 source=h3 target=subsection reason=relative-heading-rank; base=h2

\subsection{第七条教规}

从异端归正者，即诺瓦提安派、福提努斯派和十四日派，无论他们曾是其中的慕道者还是已领受共融者，除非他们诅咒一切异端，尤其是他们所持守的异端，否则不得被接纳；此后，他们中曾被称为共融者的人，在彻底学习信仰标记并以圣油傅油后，方可领受神圣奥迹。\par

我保留了文本中的\Quote{福提努斯派}一词，尽管它是否为插入语尚不确定。他们确实在圣三论上持有异端观点，因此与教规中提到的其他持正统观点的持不同意见者不同。还值得注意的是，该词未出现在\Person{费朗杜斯}的《教规摘要》（\WorkTitle{教规摘要}，第177条）中，也未出现在\Person{伊西多尔}的译本中。此外，在卢卡的一份拉丁文手稿以及巴黎的一份手稿中（如\Person{曼西}所注，卷五，第585页；卷二，第591页），均缺少该词。\Person{巴罗尼乌斯}、\Person{比尼乌斯}和\Person{雷米·塞利耶}均拒绝接受它。\par

\Quote{慕道者}一词在许多希腊手稿中缺失，但出现在\Person{巴尔萨蒙}的著作中，而且\Person{狄奥尼修斯}和\Person{伊西多尔}的文本中也有。\par

% structure: p0027 source=h3 target=subsection reason=relative-heading-rank; base=h2

\subsection{第八条教规}

从所谓的弗里吉亚派异端归正者，即使他们被该派视为神职人员，甚至被称为最高领袖，也应由教会的主教和司铎谨慎地加以教导并为他们施洗。\par

赫费勒：\Quote{本主教会议在此宣布孟他努派的圣洗无效，而在前一条法规中，它承认诺瓦提安派和十四日派的圣洗为有效。由此可见，孟他努派在有关天主圣三的教义上被怀疑有异端。然而，古代教会的其他一些权威人士则有不同判断，长期以来，教会中一直存在一个问题，即是否应将孟他努派的圣洗视为有效。亚历山大的狄约尼削大帝赞成其有效性；但本主教会议和第二届大公会议拒绝承认其为有效，更不必说公元235年的依科尼雍主教会议宣布所有异端者的圣洗无效。古代教会之所以存在这种不确定性，原因如下：（一）一方面，孟他努派，尤其戴尔都良，声称他们与天主教徒持有相同的信仰和圣事，特别是相同的圣洗（\OriginalTerm{相同的圣洗圣事}{eadem lavacri sacramenta}）。圣厄皮法尼乌斯赞同这一点，并证明孟他努派在圣父、圣子、圣神方面的教导与天主教会相同。（二）然而，其他教父对他们评价较低，原因是孟他努派常以如此模棱两可的方式表达自己，以至于他们可能——不，必须——被认为完全将圣神与孟他努等同起来。因此，戴尔都良在引用孟他努的言论时，实际上说：\InnerQuote{护慰者说话}；因此，费尔米里安、耶路撒冷的济利禄、巴西略大帝及其他教父确实以此等同来指责孟他努派，并因此认为他们的圣洗无效。（三）巴西略大帝在这方面走得最远，他坚持认为孟他努派是因圣父、圣子及孟他努和普黎斯齐拉之名施予圣洗。但正如蒂尔蒙所推测，巴西略很可能只是基于他假设孟他努派将孟他努等同于圣神，而编造出这些关于他们圣洗方式的奇怪故事；并且，正如巴罗尼乌斯所主张，同样\InnerQuote{可能的是，孟他努派并未改变圣洗的形式。然而，即使承认这一切，他们关于孟他努和圣神的模棱两可的表述，本身就足以使宣告他们的圣洗无效成为明智之举。}（四）此外，相当多的孟他努派，即埃斯基内斯学派，陷入了撒伯里乌主义，因此他们的圣洗明确无效。（参见韦策与韦尔特\WorkTitle{教会词典}中\InnerQuote{孟他努}词条；本人（即赫费勒）撰文）。}\par

% structure: p0030 source=h3 target=subsection reason=relative-heading-rank; base=h2

\subsection{第九条}

教会成员不得在墓地聚会，也不得参加任何异端者的所谓殉道者纪念处的祈祷或礼仪；若有如此行为者，若是享有共融者，应暂予绝罚；但若他们悔改并承认自己犯了罪，则可被接纳。\par

\EditorialNote{巴尔萨蒙：\Quote{正如第六条法规禁止异端者进入天主的殿宇，本条法规禁止信友前往异端者的墓地，这些墓地被他们称为殉道者纪念处……因为在迫害时期，某些异端者自称基督徒，甚至殉道至死，因此那些赞同他们意见的人称他们为殉道者。}}\par

% structure: p0033 source=h3 target=subsection reason=relative-heading-rank; base=h2

\subsection{第十条}

教会成员不得不加区别地将自己的子女嫁给异端者。\par

% structure: p0035 source=h3 target=subsection reason=relative-heading-rank; base=h2

\subsection{第十一条}

被称为\OriginalTerm{长老妇}{Presbytides}或女主席者，不得在教会中设立。\par

巴尔萨蒙：\par

在古代，某些可敬的妇女（\OriginalTerm{长老妇}{\GreekText{πρεσβύτιδες}}）坐在天主教教堂中，负责监督其他妇女保持良好和得体的秩序。但由于她们滥用本应正确使用的事物，或因傲慢或卑鄙的利己之心，产生了丑闻。因此，教父们禁止此后在教会中存在任何被称为长老妇或主席的妇女。为避免有人反对说，在女隐修院中必须有一位妇女管理其余人，应记住，她们对天主的自我舍弃和削发礼使她们虽人数众多，却被视为一个身体；她们的一切只关乎灵魂的得救。然而，妇女在天主教教堂中教导，那里聚集着众多男子和不同意见的妇女，是极不体面且有害的。\par

赫费勒：\par

赫费勒：此处意图不明，本条法规有不同的解释。首先，\OriginalTerm{长老妇}{\GreekText{πρεσβύτιδες}}和\OriginalTerm{主席}{\GreekText{προκαθήμεναι}}（\Quote{长老妇}和女主席）是什么意思？我认为厄皮法尼乌斯在他反对科利里迪安派的论文（\WorkTitle{异端}第七十九卷第四节）中首先阐明了这一点，他说：\Quote{妇女从未被允许献祭，正如科利里迪安派擅自所为，只允许她们服务。因此，教会中只有女执事，即使其中最年长的被称为\InnerQuote{长老妇}，这个词也必须与女司铎明确区分。后者意为女祭司（\OriginalTerm{女祭司}{\GreekText{ἱερίσσας}}），而\InnerQuote{长老妇}仅指她们的年龄，如同长者。}据此，本条法规似乎涉及作为其他女执事监督者（\OriginalTerm{主席}{\GreekText{προκαθήμεναι}}）的高级女执事；随后的文字可能意味着今后不再设立此类高级女执事或女长老，可能是因为她们常常越权。\par

然而，内安德、富克斯等人认为，更可能的是，本条法规中的这些术语仅指女执事，因为在教会中她们也习惯于管理女性信徒（因此得名\Quote{主席}）；并且，按照圣保禄的规定，只有年过六十的寡妇才应被选担任此职（因此称为\Quote{长老妇}）。我们可以补充说，宗徒的这一指示后来并未被严格遵守，但仍反复强调应只选年长者作为女执事。例如，加采东大公会议在其第十五条法规中要求女执事至少年满四十岁，而德奥多西皇帝甚至规定六十岁。\par

假定此条令仅仅涉及女执事，那么便产生了一个新的疑问：最后几句话——\Quote{她们不应在教会中受职}应如何理解。因为它可能意味着\Quote{从今以后不得再任命女执事}；或者\Quote{今后她们不应再在教会中受隆重祝圣}。然而，第一种解释将与希腊教会在老底嘉会议后很久仍有女执事的事实相矛盾。例如，692年的\Emph{特鲁洛}会议（第十四条）规定\Quote{年龄未满四十岁者不得被祝圣为女执事}。因此，第二种解释\Quote{她们不应在教会中受隆重祝圣}似乎更合理，而尼安德也明确倾向于这种看法。确实，后来的几次会议明确禁止了旧日对女执事授予某种祝圣的做法，例如第一次\Emph{奥朗日}会议（441年，第二十六条）的措辞——\Quote{diaconæ omnimodis non ordinandæ}；同样还有517年的埃帕翁会议（第二十一条）和533年的第二次奥尔良会议（第十八条）。但在希腊教会中，至少直至特鲁洛会议（第十四条）仍有祝圣（\OriginalTerm{被祝圣}{\GreekText{χειροτονεῖσθαι}}）发生。然而，老底嘉的这条令并未提及隆重奉献，更非祝圣，而仅涉及\OriginalTerm{任命}{\GreekText{καθίστασθαι}}。这些理由促使我们回到对本条令的第一种解释，即理解为其禁止此后任命任何首席女执事或\Quote{\OriginalTerm{首席女执事}{presbytides}}。\par

佐纳拉斯与巴尔萨蒙给出了另一种解释。依他们之见，这些\Quote{\OriginalTerm{首席女执事}{presbytides}}并非首席女执事，而是一般的老妇人（\Emph{来自民众}），被赋予在教会中监督女性的职责。然而，老底嘉会议废除了这一安排，大概是因为她们滥用其职位，出于骄傲、牟利、受贿等目的。\par

% structure: p0044 source=h3 target=subsection reason=relative-heading-rank; base=h2

\subsection{第十二条}

主教应由都主教及邻近主教之判断，在对其信德根基与正直生活之操守经过长久考验后，任命管理教会事务。\par

% structure: p0046 source=h3 target=subsection reason=relative-heading-rank; base=h2

\subsection{第十三条}

将被任命为司铎者的选举，不得委诸群众。\par

% structure: p0048 source=h3 target=subsection reason=relative-heading-rank; base=h2

\subsection{第十四条}

在逾越节庆日，不得以\OriginalTerm{祝福礼}{eulogiæ}的方式将圣物送往其他教区。\par

% structure: p0050 source=h3 target=subsection reason=relative-heading-rank; base=h2

\subsection{第十五条}

除正规歌咏者外，任何人不得在教会中歌唱；正规歌咏者应登上读经台，按歌本歌唱。\par

% structure: p0052 source=h3 target=subsection reason=relative-heading-rank; base=h2

\subsection{第十六条}

在安息日（即星期六），应诵读福音书及其他圣经。\par

在教会圣咏安排尚未确定之前，安息日本不惯常诵读福音书或其他圣经。但基于禁止在安息日斋戒或跪拜的条令，当时并无礼仪，以便尽量欢庆。教父们禁止此做法，并规定在安息日应举行完整的日课。\par

尼安德（\Emph{教会史}，第二版，第三卷，第565页及以下）除上述解释外，还提出了另一种解释，即：古代教会许多地方将每个星期六作为纪念创造的节日来庆祝。尼安德还推测，可能有些犹太化者只在安息日诵读旧约；但他本人也指出，在这种情况下，\OriginalTerm{福音（复数）}{\GreekText{εὐαγγέλια}}与\OriginalTerm{其他圣经书卷}{\GreekText{ἑτέρων} \GreekText{γραφῶν}}会需要冠词。\par

\Signature{Van Espen.}\par

在希腊人中，安息日的庆祝方式与主日完全相同，只是涉及停工方面有所区别，因此会议希望，如同主日一样，在其余读经之后，应接着诵读福音。\par

因为显然，教会的意图是整个日课都旨在启迪和教育信众，尤其在节日，信众往往大量聚集。\par

在此我们可以注意我们当前（西方）礼规的起源：在主日和节日，通常在时辰礼仪中与其余圣经一同诵读福音，而在平日或平日及\Quote{简单}庆日的礼规中则不然。\par

% structure: p0060 source=h3 target=subsection reason=relative-heading-rank; base=h2

\subsection{第十七条}

在会众中，不得连续诵唱圣咏，每首圣咏之后应插入一段读经。\par

% structure: p0062 source=h3 target=subsection reason=relative-heading-rank; base=h2

\subsection{第十八条}

无论是在第九时辰祈祷还是在晚祷，应始终诵念相同的祈祷经文。\par

赫费勒：\Quote{有些节庆在第九时辰结束，另一些仅在傍晚结束，两者皆以祈祷结束。此处会议规定在两种情况下应使用相同的祈祷。范埃斯彭如此解释文本的措辞，我认为正确。但希腊评注者认为会议命令在所有地方使用相同的祈祷，从而排除个人的任意性。据此，一致性规则涉及地域；而依范埃斯彭，第九时辰祈祷与晚祷应是相同的。然而，若此解释正确，会议不应仅提及第九时辰祈祷与晚祷的祈祷，而应笼统地说\InnerQuote{所有教区应使用相同形式的祈祷}。}\par

% structure: p0065 source=h3 target=subsection reason=relative-heading-rank; base=h2

\subsection{第十九条}

在主教的讲道之后，应首先单独为慕道者祈祷；慕道者退出后，为那些处于补赎中的人祈祷；这些人经过（主教的）覆手并离去后，随后应献上三篇信友祷文，第一篇完全默念，第二、第三篇高声诵念，然后行平安礼。司铎向主教行平安礼后，平信徒再彼此行礼，然后完成神圣祭献。唯有圣职人员方可前往祭台并（在此）领受共融。\par

% structure: p0067 source=h3 target=subsection reason=relative-heading-rank; base=h2

\subsection{第二十条}

执事在司铎面前坐下是不合宜的，除非司铎命其坐下。同样，执事应受到副执事及全体（下级）圣职人员的尊敬。\par

% structure: p0069 source=h3 target=subsection reason=relative-heading-rank; base=h2

\subsection{第二十一条}

副执事无权进入\OriginalTerm{圣器室}{Diaconicum}，亦不得触摸主的器皿。\par

\EditorialNote{赫费勒：对于\Quote{圣器室}此处是指执事在礼仪中站立之处，还是通常所称的圣器室（相当于今日的祭衣房），存在疑问。在这圣器室内保存着圣器与祭衣；由于本条令的最后部分特别提到这些，我毫不怀疑\Quote{圣器室}必须指祭衣房。此外，本条令不过是对该规则的具体表述，即副执事不得擅行执事的职务。}\par

% structure: p0072 source=h3 target=subsection reason=relative-heading-rank; base=h2

\subsection{第二十二条}

副执事无权佩戴\OriginalTerm{披肩}{orarium}，也不得离开门口。\par

\EditorialNote{据早年，按\Person{Zonaras}和\Person{Balsamon}所述，副执事的职责是站在教堂门口，并在礼仪的适当时刻领入和领出望教者和补赎者。\Person{Zonaras}评论说，若此习俗如同许多其他古老习俗一样完全改变并废弃，无人须惊讶。}\par

% structure: p0075 source=h3 target=subsection reason=relative-heading-rank; base=h2

\subsection{第二十三条例}

读者和歌咏者无权佩戴\OriginalTerm{披肩}{orarium}，并以此装束诵读或歌咏。\par

% structure: p0077 source=h3 target=subsection reason=relative-heading-rank; base=h2

\subsection{第二十四条例}

司祭职中的任何人，从司铎到执事，以及教会等级中依次的副执事、读者、歌咏者、驱魔者、守门者，或任何\OriginalTerm{精修者}{Ascetics}，都不得进入酒馆。\par

% structure: p0079 source=h3 target=subsection reason=relative-heading-rank; base=h2

\subsection{第二十五条例}

副执事不得授予饼，也不得祝福杯。\par

\EditorialNote{根据\Person{Hefele}：\Quote{按照《宗徒宪章》，共融礼以如下方式施行：主教将圣饼交给每人，说：主的身体，领受者答：阿们。然后执事递过圣爵，说：基督的血，生命的爵，领受者再答：阿们。此递圣爵时所说的\Quote{基督的血}等，在劳底嘉规例中被称为祝福（\OriginalTerm{祝福}{\GreekText{εὐλογεῖν}}）。}}。希腊注释家\Person{Aristenus}按照此，而且非常正确地，给出了本条规例的含义。\par

% structure: p0082 source=h3 target=subsection reason=relative-heading-rank; base=h2

\subsection{第二十六条例}

凡未经主教晋升者，无论在教堂或私宅，均不得施行驱魔礼。\par

% structure: p0084 source=h3 target=subsection reason=relative-heading-rank; base=h2

\subsection{第二十七条例}

受邀参加爱筵者，无论是司祭人员、神职人员或平信徒，均不得带走自己的份额，因为这是对教会秩序的羞辱。\par

\Signature{Hefele}\par

\Person{Van Espen}翻译为：\Quote{在教会中担任任何职务者，无论是神职人员或平信徒，}并援引事实，即早在早期，希腊人中许多人未经祝圣就在教会中任职，如同现今我们的圣器管理员和辅祭一样。但我不同意\Person{Van Espen}认为\Quote{司祭人员}泛指教会中任何任职者，而仅指较高等级的神职人员，即司铎和执事，正如前第二十四条例中，只明确将司铎和执事列于\OriginalTerm{司祭}{\GreekText{ἱερατικοῖς}}中，并与其他（较低级）神职人员区分。后来在第三十条规例中，也类似提到三个不同等级：\OriginalTerm{司祭}{\GreekText{ἱερατικοί}}、\OriginalTerm{神职人员}{\GreekText{κληρικοί}}和\OriginalTerm{精修者}{\GreekText{ἀσκηταί}}。\par

将\OriginalTerm{爱筵}{agape}的剩余物带走在此被禁止，因为一方面显示贪婪，另一方面或许被视为亵渎。\par

% structure: p0089 source=h3 target=subsection reason=relative-heading-rank; base=h2

\subsection{第二十八条例}

不允许在主的殿宇或教堂内举行所谓爱筵，也不得在天主屋内进食和铺床。\par

% structure: p0091 source=h3 target=subsection reason=relative-heading-rank; base=h2

\subsection{第二十九条例}

基督徒不得因安息日休息而犹太化，而应在该日工作，反而应尊崇主日；若可能，则在那日作为基督徒休息。但若有人被发现是犹太化者，愿他从基督被开除教籍。\par

% structure: p0093 source=h3 target=subsection reason=relative-heading-rank; base=h2

\subsection{第三十条规例}

任何司祭、神职人员、精修者，或任何基督徒、平信徒，都不可与妇女同浴；因为这在异教徒中是最可耻的。\par

% structure: p0095 source=h3 target=subsection reason=relative-heading-rank; base=h2

\subsection{第三十一条规例}

与各类异端者通婚是不合法的，也不可将我们的儿女嫁娶给他们；但若他们承诺成为基督徒，则可从他们中娶嫁。\par

% structure: p0097 source=h3 target=subsection reason=relative-heading-rank; base=h2

\subsection{第三十二条规例}

接受异端者的\OriginalTerm{祝福}{eulogiæ}是不合法的，因为它们与其说是\OriginalTerm{祝福}{eulogiæ}，不如说是\OriginalTerm{愚妄}{\GreekText{ἀλογίαι}}。\par

% structure: p0099 source=h3 target=subsection reason=relative-heading-rank; base=h2

\subsection{第三十三条规例}

任何人不得与异端者或分裂者一同祈祷。\par

% structure: p0101 source=h3 target=subsection reason=relative-heading-rank; base=h2

\subsection{第三十四条规例}

任何基督徒不可离弃基督的殉道者，而转向假殉道者，即那些属于异端者的，或那些先前曾是异端者的；因为他们与天主疏远。因此，那些追随他们的人，应被开除教籍。\par

\Signature{Hefele}\par

本条规例禁止尊崇不属于正统教会的殉道者。弗里吉亚的孟他努派殉道者人数众多，可能是本条规例的起因。\par

我翻译的\Quote{那些先前曾是异端者}这一短语给所有译者带来很大困难，几乎没有两人意见一致。\Person{Hammond}读作\Quote{那些曾被认作异端者的人}，\Person{Fulton}同意他，但错误地（如我所认为）省略了\Quote{to}。\Person{Lambert}翻译为\Quote{那些从前是异端者的人}，并且正确。\Person{Van Espen}同意他，即\OriginalTerm{那些先前曾是异端者}{vel eos qui prius heretici fuere}。\par

% structure: p0106 source=h3 target=subsection reason=relative-heading-rank; base=h2

\subsection{第三十五条规例}

基督徒不可离弃天主的教会，并去呼求天使、聚集集会，这些事被禁止。因此，若有人被发现从事这种隐秘的偶像崇拜，应被开除教籍；因为他已离弃了我们的主耶稣基督、天主子，而转向了偶像崇拜。\par

应注意，某些极具权威和古老的拉丁译本中\OriginalTerm{角落}{angulos}代替了\OriginalTerm{天使}{angelos}。这指在角落、隐藏地、秘密地进行这些偶像崇拜仪式，如拉丁文\OriginalTerm{秘密地}{occulte}。但此读法虽在拉丁文中如此受尊重，却无希腊文权威支持。\par

本条规例常被在争论中用来谴责天主教会始终给予天使的敬礼，但那些欲如此使用本规例的人，应解释这些解释如何能与同一次大公会议明显批准的殉道者敬礼相符合；以及若本条规例谴责了当时东西方教会的普遍实践，它又如何能被尼西亚第二次大公会议的教父们所接受。参阅\Person{Forbes}，\WorkTitle{Considerationes Modestæ}。\par

% structure: p0110 source=h3 target=subsection reason=relative-heading-rank; base=h2

\subsection{第三十六条规例}

属司祭或神职人员者，不得为巫术师、行邪术者、占星家或占星术士；也不得制作所谓护身符，那是对自己灵魂的枷锁。凡佩戴此类者，我们命令将他们逐出教会。\par

% structure: p0112 source=h3 target=subsection reason=relative-heading-rank; base=h2

\subsection{第三十七条规例}

接受从犹太人或异端者筵席送来的份额是不合法的，也不可与他们一同宴饮。\par

% structure: p0114 source=h3 target=subsection reason=relative-heading-rank; base=h2

\subsection{第三十八条规例}

接受犹太人的无酵饼是不合法的，也不可参与他们的不敬。\par

% structure: p0116 source=h3 target=subsection reason=relative-heading-rank; base=h2

\subsection{第三十九条规例}

与异教徒一同宴饮是不合法的，也不可参与他们的无神。\par

% structure: p0118 source=h3 target=subsection reason=relative-heading-rank; base=h2

\subsection{第四十条规例}

被召赴大公会议的主教不得犯藐视罪，而应出席，或教授，或受教，以促进教会及他人的革新。若有人犯了藐视罪，他将定自己的罪，除非因健康欠佳被阻延。\par

\Signature{Hefele}\par

通常\OriginalTerm{疾病}{\GreekText{ἀνωμαλία}}被理解为疾病，\Person{Dionysius Exiguus}和\Person{Isidore}翻译为前者\OriginalTerm{疾病}{ægritudinem}，后者\OriginalTerm{虚弱}{infirmitatem}。但\Person{Balsamon}公正地评论说，该术语有更广泛的意义，除疾病情况外还包括其他不可避免的阻碍或障碍。\par

% structure: p0122 source=h3 target=subsection reason=relative-heading-rank; base=h2

\subsection{第四十一条规例}

无主教之命，任何司祭或神职人员不得出行。\par

% structure: p0124 source=h3 target=subsection reason=relative-heading-rank; base=h2

\subsection{第四十二条规例}

任何司铎或圣职人员，若无教会公函，不得旅行。\par

% structure: p0126 source=h3 target=subsection reason=relative-heading-rank; base=h2

\subsection{第四十三条例}

副执事不得离开门口去参与祈祷，即使短暂时间也不可。\par

% structure: p0128 source=h3 target=subsection reason=relative-heading-rank; base=h2

\subsection{第四十四条例}

妇女不得走近祭台。\par

% structure: p0130 source=h3 target=subsection reason=relative-heading-rank; base=h2

\subsection{第四十五条例}

［领洗候选人］在四旬期第二周之后，不得接受。\par

\Person{Van Espen}。\par

为理解本条条例，必须记得：那些愿意成为天主教徒并受洗的外邦人，起初由传道员私下教导。此后，在获得一些基督宗教知识后，他们被接纳参与主教在教会内公开讲授，因而被称为\OriginalTerm{听讲者}{Audientes}，并首次确切地被称为\OriginalTerm{慕道者}{Catechumens}。但当这些慕道者在此等级中经过足够时间并经受考验后，他们被允许升至更高等级，称为\OriginalTerm{跪拜者}{Genuflectentes}。\par

当他们在此顺序中完成操练后，负责他们的传道员便将他们带到主教前，为在圣安息日（即复活节前夕）领受洗礼；同时慕道者呈上他们的名字，以便被登记在即将来临的圣安息日受洗。\par

此外，我们从圣奥斯定（讲道集第十三篇，《致新领洗者》）得知，呈递名字的时间是四旬期的开始。\par

因此，本公会议借此条例规定：凡未在四旬期开始、而是在两周过后才呈递名字者，不得在下一个圣安息日获准受洗。\par

% structure: p0137 source=h3 target=subsection reason=relative-heading-rank; base=h2

\subsection{第四十六条例}

即将受洗者必须背诵信仰（信经），并在一周的第五日向主教或司铎复述。\par

\Person{Hefele}。\par

怀疑本文中所谓\Quote{星期四}是指圣周星期四，还是慕道者接受教导期间每个星期四。希腊注释者支持后者，但狄奥尼修斯·埃克西古斯与伊西多尔，以及其后的宾厄姆，则支持前者，且很可能正确。本条例后被特鲁罗会议在其第七十八条条例中重申。\par

% structure: p0141 source=h3 target=subsection reason=relative-heading-rank; base=h2

\subsection{第四十七条例}

患病时受洗而后痊愈者，必须背诵信经，并明白神圣恩赐已赐予他们。\par

% structure: p0143 source=h3 target=subsection reason=relative-heading-rank; base=h2

\subsection{第四十八条例}

受洗者必须在洗礼后以天上之圣化圣油（\OriginalTerm{圣化圣油}{chrism}）傅油，并分享基督的王国。\par

% structure: p0145 source=h3 target=subsection reason=relative-heading-rank; base=h2

\subsection{第四十九条例}

在四旬期内，除安息日和主日外，不得献饼。\par

\Person{Hefele}。\par

本条例被特鲁罗会议第五十二条条例重申，规定在四旬期的平日，只应举行\OriginalTerm{预先圣化弥撒}{Missa Præsanctificatorum}，正如希腊人在所有补赎和哀悼日仍然保持的习惯，因为他们认为在这些日子不适于举行完整的礼仪，正如利奥·阿拉提乌斯所说，这是因为祝圣是喜乐的行为。然而，与上述第十六条条例比较，表明星期六是特别的例外。\par

在赫费勒所提到的星期六和星期日之外，还必须加上圣母领报瞻礼，该瞻礼始终以完整的礼仪隆重举行，即使落在圣周五也不例外。\par

% structure: p0150 source=h3 target=subsection reason=relative-heading-rank; base=h2

\subsection{第五十条例}

在四旬期最后一周的第五日（即濯足星期四）不可破斋，以免辱没整个四旬期；但必须在整个四旬期内守斋，只可吃干食。\par

% structure: p0152 source=h3 target=subsection reason=relative-heading-rank; base=h2

\subsection{第五十一条例}

在四旬期内不得庆祝殉道者诞辰，但应在安息日和主日纪念圣殉道者。\par

% structure: p0154 source=h3 target=subsection reason=relative-heading-rank; base=h2

\subsection{第五十二条例}

婚姻和生日宴会不得在四旬期内举行。\par

\Person{Hefele}。\par

在本条例中，\Quote{生日宴会}不应理解为前一条例中的\OriginalTerm{殉道者诞辰}{natalitia martyrum}，而是君王的生日宴会。本条与前一条例一起，在第六世纪由布拉卡拉（今葡萄牙布拉加）主教马丁重新颁布。\par

% structure: p0158 source=h3 target=subsection reason=relative-heading-rank; base=h2

\subsection{第五十三条例}

\Emph{基督徒}参加婚礼时，不得参与放荡的舞蹈，而应端庄地用餐或进早餐，合乎基督徒的身份。\par

% structure: p0160 source=h3 target=subsection reason=relative-heading-rank; base=h2

\subsection{第五十四条例}

司铎和圣职人员不得观看婚礼或宴席上的戏剧；但在演员入场之前，他们必须起身离开。\par

\Person{Aristenus}。\par

基督徒被劝戒，参加婚礼时应端庄地宴饮，不跳舞，也不\OriginalTerm{拍手}{\GreekText{βαλλίζειν}}，即拍手并发出响声。因为这与基督徒的身份不相称。但献身者不得观看婚礼上的戏剧，在\OriginalTerm{戏剧演员}{thymelici}开始之前，他们必须出去。\par

% structure: p0164 source=h3 target=subsection reason=relative-heading-rank; base=h2

\subsection{第五十五条例}

司铎、圣职人员，甚至平信徒，均不得凑份子举办饮酒聚会。\par

这些餐食的费用由多人凑份子共同支付，伊西多尔称之为\Quote{symbola}（\OriginalTerm{聚餐}{symbola}），而梅利努斯和克拉贝称之为\Quote{comissalia}（\OriginalTerm{聚餐}{comissalia}），尽管更常见的形式是\Quote{commensalia}（\OriginalTerm{聚餐}{commensalia}）或\Quote{comessalia}（\OriginalTerm{聚餐}{comessalia}）。参见杜康吉《词汇集》，词条\Emph{Commensalia}和\Emph{Confertum}。\par

% structure: p0167 source=h3 target=subsection reason=relative-heading-rank; base=h2

\subsection{第五十六条例}

司铎不得在主教进入\OriginalTerm{至圣所}{bema}之前先行进入并就座；他们必须与主教一同进入，除非主教因病留在家中或不在。\par

若不歪曲其含义，很难翻译本条条例。它并非规定教会游行中的尊卑次序，而是完全不同的内容，即：当主教进入圣所时，他不应独自一人走进一个已被占用的地方，而是应带着圣职人员作为荣誉护卫。他们应走在他前面还是后面，只是地方习俗的问题，原则\OriginalTerm{年幼者在前}{juniores priores}并未普遍通行。\par

% structure: p0170 source=h3 target=subsection reason=relative-heading-rank; base=h2

\subsection{第五十七条例}

\Emph{主教}不得在乡村或偏远地区任命，但可任命巡视员；已经任命者，未经城市主教同意，不得行事。同样，司铎未经主教同意，不得行事。\par

% structure: p0172 source=h3 target=subsection reason=relative-heading-rank; base=h2

\subsection{第五十八条例}

主教或司铎不得在任何私人房屋中献\OriginalTerm{祭献}{Oblation}。\par

% structure: p0174 source=h3 target=subsection reason=relative-heading-rank; base=h2

\subsection{第五十九条例}

凡由私人编撰的诗篇及任何未列入正典的书，均不得在教会中诵读，只可诵读新旧约的正典经书。\par

% structure: p0176 source=h3 target=subsection reason=relative-heading-rank; base=h2

\subsection{教规六十}

\EditorialNote{[\Emph{注——本教规的真实性最为可疑。}]}\par

以下是规定诵读的旧约全部经书：1. \WorkTitle{创世纪}；2. \WorkTitle{出谷纪}；3. \WorkTitle{肋未纪}；4. \WorkTitle{户籍纪}；5. \WorkTitle{申命纪}；6. \WorkTitle{若苏厄书}；7. \WorkTitle{民长纪}、\WorkTitle{卢德传}；8. \WorkTitle{艾斯德尔传}；9. \WorkTitle{列王纪第一}、\WorkTitle{列王纪第二}；10. \WorkTitle{列王纪第三}、\WorkTitle{列王纪第四}；11. \WorkTitle{编年纪上}、\WorkTitle{编年纪下}；12. \WorkTitle{厄斯德拉上}、\WorkTitle{厄斯德拉下}；13. \WorkTitle{圣咏集}；14. \WorkTitle{箴言}；15. \WorkTitle{训道篇}；16. \WorkTitle{雅歌}；17. \WorkTitle{约伯传}；18. \WorkTitle{十二小先知书}；19. \WorkTitle{依撒意亚}；20. \WorkTitle{耶肋米亚}、\WorkTitle{巴路克}、\WorkTitle{哀歌}、\WorkTitle{耶肋米亚书信}；21. \WorkTitle{厄则克耳}；22. \WorkTitle{达尼尔}。\par

以下是新约的经书：1. \WorkTitle{玛窦福音}、\WorkTitle{马尔谷福音}、\WorkTitle{路加福音}、\WorkTitle{若望福音}；2. \WorkTitle{宗徒大事录}；3. 七封公函：\WorkTitle{雅各伯书}、\WorkTitle{伯多禄前书}、\WorkTitle{伯多禄后书}、\WorkTitle{若望一书}、\WorkTitle{若望二书}、\WorkTitle{若望三书}、\WorkTitle{犹达书}；4. 十四封保禄书信：\WorkTitle{罗马人书}、\WorkTitle{格林多前书}、\WorkTitle{格林多后书}、\WorkTitle{迦拉达书}、\WorkTitle{厄弗所书}、\WorkTitle{斐理伯书}、\WorkTitle{哥罗森书}、\WorkTitle{得撒洛尼前书}、\WorkTitle{得撒洛尼后书}、\WorkTitle{希伯来书}、\WorkTitle{弟茂德前书}、\WorkTitle{弟茂德后书}、\WorkTitle{弟铎书}、\WorkTitle{费肋孟书}。\par

\SourceRecord{Translated by Henry Percival.；Nicene and Post-Nicene Fathers, Second Series；Vol. 14.；Edited by Philip Schaff and Henry Wace.；Buffalo, NY: Christian Literature Publishing Co.,；1900.；<http://www.newadvent.org/fathers/3806.htm>.}
